% --- Archon physics formalization source begin ---
% archon:physics
% archon:covers PhyXMiniProblems/problem_phyx_mini_0050.lean
% archon:source-report reports/phyx_mini/problem_phyx_mini_0050.source.json

\chapter{Physics problem phyx\_mini\_0050}
\label{ch:PhyXMiniProblems_problem_phyx_mini_0050}

\paragraph{Problem source.}
\#\# Physical scenario

The figure shows a laser beam deflected by a 30\textasciicircum{}\textbackslash{}circ\textbackslash{}text\{-\}60\textasciicircum{}\textbackslash{}circ\textbackslash{}text\{-\}90\textasciicircum{}\textbackslash{}circ prism.

\#\# Question

What is the prism’s index of refraction?

\#\# Answer choices

- A: A: 0.64

- B: B: 1.59

- C: C: 0.89

- D: D: 1.24

\#\# Figure caption (auxiliary; use the image as primary evidence)

This image depicts a geometric illustration involving a laser beam and a triangular prism. Here's a detailed description:

- The **laser beam** is represented by a horizontal red rectangle on the left side of the image.
- The beam passes through a **triangular prism** shaded in blue. This prism has three angles labeled: 
  - Top angle: 30°
  - Bottom angle: 60°
- As the laser exits the prism, its path bends, creating an angle of 22.6° with the original path.
- An arrow at the end of the beam's path indicates the direction of the refracted beam.
- The text “Laser beam” is located next to the entering beam.

This diagram illustrates the refraction of a laser beam as it passes through a prism, highlighting the angular relationships.

\#\# Dataset metadata

Geometrical Optics; Physical Model Grounding Reasoning, Spatial Relation Reasoning

\paragraph{Recorded answer/context.}
B: B: 1.59

\paragraph{Figure/image path.}
/root/proposal\_for\_physic/science-mango/phyx\_mini\_run/phyx\_data/test\_image/50.png

\paragraph{Formalization target.}
create a compiling Lean file with sorry bodies at `PhyXMiniProblems/problem\_phyx\_mini\_0050.lean`. The Lean declarations must preserve the physical quantities, dimensions or dimensional roles, figure labels, governing-law hypotheses, and final relation expressed by this problem.
Use Mathlib/Physlib names found through LeanExplore where available. If a physics API is missing, introduce faithful local abstractions rather than scalar placeholder aliases.

\begin{theorem}[Physics formalization target]
\label{thm:physics:phyx_mini_0050:target}
\lean{PhyXMiniProblems.ProblemPhyXMini0050.problem_phyx_mini_0050}
\uses{def:physics:phyx-mini-0050:phyxminiproblems-problemphyxmini0050-dimensionlessreadout, def:physics:phyx-mini-0050:phyxminiproblems-problemphyxmini0050-radiansofdegrees, def:physics:phyx-mini-0050:phyxminiproblems-problemphyxmini0050-prismlasersetup, def:physics:phyx-mini-0050:phyxminiproblems-problemphyxmini0050-matchesfigurereadouts, def:physics:phyx-mini-0050:phyxminiproblems-problemphyxmini0050-hasdepictedprismgeometry, def:physics:phyx-mini-0050:phyxminiproblems-problemphyxmini0050-hasdepictedraypath, def:physics:phyx-mini-0050:phyxminiproblems-problemphyxmini0050-hasphysicalrefractiveindices, def:physics:phyx-mini-0050:phyxminiproblems-problemphyxmini0050-satisfiessnelllawatprismfaces, def:physics:phyx-mini-0050:phyxminiproblems-problemphyxmini0050-matchesanswertonearesthundredth}
The assigned autoformalize agent should translate this physics problem into Lean declarations in the covered file, with theorem and lemma proof bodies written as `by sorry`.
\end{theorem}
\begin{proof}
This is an autoformalization task, not a proof task. Produce statements that can later be proved by the physics prover without weakening the physical model.
\end{proof}

% --- Archon declaration topology begin ---
\section{Declaration topology}
\begin{definition}[Dimensionless Quantity]
  \label{def:physics:phyx-mini-0050:phyxminiproblems-problemphyxmini0050-dimensionlessquantity}
  \lean{PhyXMiniProblems.ProblemPhyXMini0050.DimensionlessQuantity}
  A unit-independent physical scalar, carried by Physlib's identity dimension
\end{definition}
\begin{proof}
  This is the declared model definition of Dimensionless Quantity.
\end{proof}
\begin{definition}[dimensionless Readout]
  \label{def:physics:phyx-mini-0050:phyxminiproblems-problemphyxmini0050-dimensionlessreadout}
  \lean{PhyXMiniProblems.ProblemPhyXMini0050.dimensionlessReadout}
  \uses{def:physics:phyx-mini-0050:phyxminiproblems-problemphyxmini0050-dimensionlessquantity}
  The real scalar readout of a dimensionless physical quantity
\end{definition}
\begin{proof}
  This is the declared model definition of dimensionless Readout.
\end{proof}
\begin{definition}[radians Of Degrees]
  \label{def:physics:phyx-mini-0050:phyxminiproblems-problemphyxmini0050-radiansofdegrees}
  \lean{PhyXMiniProblems.ProblemPhyXMini0050.radiansOfDegrees}
  Convert the degree labels printed in the figure to radian values
\end{definition}
\begin{proof}
  This is the declared model definition of radians Of Degrees.
\end{proof}
\begin{definition}[Optical Region]
  \label{def:physics:phyx-mini-0050:phyxminiproblems-problemphyxmini0050-opticalregion}
  \lean{PhyXMiniProblems.ProblemPhyXMini0050.OpticalRegion}
  The two physical regions crossed by the laser beam
\end{definition}
\begin{proof}
  This is the declared model definition of Optical Region.
\end{proof}
\begin{definition}[Prism Vertex]
  \label{def:physics:phyx-mini-0050:phyxminiproblems-problemphyxmini0050-prismvertex}
  \lean{PhyXMiniProblems.ProblemPhyXMini0050.PrismVertex}
  The three vertices of the triangular prism in the primary image
\end{definition}
\begin{proof}
  This is the declared model definition of Prism Vertex.
\end{proof}
\begin{definition}[Prism Face]
  \label{def:physics:phyx-mini-0050:phyxminiproblems-problemphyxmini0050-prismface}
  \lean{PhyXMiniProblems.ProblemPhyXMini0050.PrismFace}
  The prism interfaces crossed by the ray, in propagation order
\end{definition}
\begin{proof}
  This is the declared model definition of Prism Face.
\end{proof}
\begin{definition}[Ray Segment]
  \label{def:physics:phyx-mini-0050:phyxminiproblems-problemphyxmini0050-raysegment}
  \lean{PhyXMiniProblems.ProblemPhyXMini0050.RaySegment}
  The three directed portions of the depicted laser path
\end{definition}
\begin{proof}
  This is the declared model definition of Ray Segment.
\end{proof}
\begin{definition}[Prism Laser Setup]
  \label{def:physics:phyx-mini-0050:phyxminiproblems-problemphyxmini0050-prismlasersetup}
  \lean{PhyXMiniProblems.ProblemPhyXMini0050.PrismLaserSetup}
  \uses{def:physics:phyx-mini-0050:phyxminiproblems-problemphyxmini0050-dimensionlessquantity, def:physics:phyx-mini-0050:phyxminiproblems-problemphyxmini0050-opticalregion, def:physics:phyx-mini-0050:phyxminiproblems-problemphyxmini0050-prismvertex, def:physics:phyx-mini-0050:phyxminiproblems-problemphyxmini0050-prismface, def:physics:phyx-mini-0050:phyxminiproblems-problemphyxmini0050-raysegment}
  The physical and geometrical quantities in the prism experiment. Directions are oriented radian angles measured counterclockwise from the rightward horizontal direction. At each crossed face, propagationNormalDirectionRad chooses the normal pointing into the forward half-space of the propagating ray. This makes incidence/refraction angles the ordinary acute differences from the relevant normal
\end{definition}
\begin{proof}
  This is the declared model definition of Prism Laser Setup.
\end{proof}
\begin{definition}[angle From Propagation Normal]
  \label{def:physics:phyx-mini-0050:phyxminiproblems-problemphyxmini0050-anglefrompropagationnormal}
  \lean{PhyXMiniProblems.ProblemPhyXMini0050.angleFromPropagationNormal}
  \uses{def:physics:phyx-mini-0050:phyxminiproblems-problemphyxmini0050-prismface, def:physics:phyx-mini-0050:phyxminiproblems-problemphyxmini0050-raysegment, def:physics:phyx-mini-0050:phyxminiproblems-problemphyxmini0050-prismlasersetup}
  The nonnegative angle between a ray segment and a chosen face normal
\end{definition}
\begin{proof}
  This is the declared model definition of angle From Propagation Normal.
\end{proof}
\begin{definition}[Matches Figure Readouts]
  \label{def:physics:phyx-mini-0050:phyxminiproblems-problemphyxmini0050-matchesfigurereadouts}
  \lean{PhyXMiniProblems.ProblemPhyXMini0050.MatchesFigureReadouts}
  \uses{def:physics:phyx-mini-0050:phyxminiproblems-problemphyxmini0050-radiansofdegrees, def:physics:phyx-mini-0050:phyxminiproblems-problemphyxmini0050-prismlasersetup}
  Numerical and directional readouts visible in the primary image: the prism vertices are 30°, 90°, and 60°, and the outgoing beam is deflected 22.6° below the original horizontal laser direction
\end{definition}
\begin{proof}
  This is the declared model definition of Matches Figure Readouts.
\end{proof}
\begin{definition}[Has Depicted Prism Geometry]
  \label{def:physics:phyx-mini-0050:phyxminiproblems-problemphyxmini0050-hasdepictedprismgeometry}
  \lean{PhyXMiniProblems.ProblemPhyXMini0050.HasDepictedPrismGeometry}
  \uses{def:physics:phyx-mini-0050:phyxminiproblems-problemphyxmini0050-radiansofdegrees, def:physics:phyx-mini-0050:phyxminiproblems-problemphyxmini0050-prismlasersetup}
  Spatial relations supplied by the right-triangle prism geometry. The entry face is vertical, hence its forward normal is horizontal. Since the sloping face meets the vertical face at the top vertex, its outward forward normal has the direction of that 30° vertex angle in the pictured coordinates
\end{definition}
\begin{proof}
  This is the declared model definition of Has Depicted Prism Geometry.
\end{proof}
\begin{definition}[Has Depicted Ray Path]
  \label{def:physics:phyx-mini-0050:phyxminiproblems-problemphyxmini0050-hasdepictedraypath}
  \lean{PhyXMiniProblems.ProblemPhyXMini0050.HasDepictedRayPath}
  \uses{def:physics:phyx-mini-0050:phyxminiproblems-problemphyxmini0050-prismlasersetup}
  The depicted propagation path: normal entry produces the horizontal internal segment, and the exiting segment is clockwise from the incident segment by the positive deflection shown in the figure
\end{definition}
\begin{proof}
  This is the declared model definition of Has Depicted Ray Path.
\end{proof}
\begin{definition}[Has Physical Refractive Indices]
  \label{def:physics:phyx-mini-0050:phyxminiproblems-problemphyxmini0050-hasphysicalrefractiveindices}
  \lean{PhyXMiniProblems.ProblemPhyXMini0050.HasPhysicalRefractiveIndices}
  \uses{def:physics:phyx-mini-0050:phyxminiproblems-problemphyxmini0050-dimensionlessreadout, def:physics:phyx-mini-0050:phyxminiproblems-problemphyxmini0050-prismlasersetup}
  The idealized surrounding air has index one, and both indices are positive
\end{definition}
\begin{proof}
  This is the declared model definition of Has Physical Refractive Indices.
\end{proof}
\begin{definition}[Satisfies Snell Law At Prism Faces]
  \label{def:physics:phyx-mini-0050:phyxminiproblems-problemphyxmini0050-satisfiessnelllawatprismfaces}
  \lean{PhyXMiniProblems.ProblemPhyXMini0050.SatisfiesSnellLawAtPrismFaces}
  \uses{def:physics:phyx-mini-0050:phyxminiproblems-problemphyxmini0050-dimensionlessreadout, def:physics:phyx-mini-0050:phyxminiproblems-problemphyxmini0050-prismlasersetup, def:physics:phyx-mini-0050:phyxminiproblems-problemphyxmini0050-anglefrompropagationnormal}
  Snell's law at the entry and exit interfaces. Every angle is measured from the local propagation normal, and every refractive index is dimensionless
\end{definition}
\begin{proof}
  This is the declared model definition of Satisfies Snell Law At Prism Faces.
\end{proof}
\begin{definition}[Answer Choice]
  \label{def:physics:phyx-mini-0050:phyxminiproblems-problemphyxmini0050-answerchoice}
  \lean{PhyXMiniProblems.ProblemPhyXMini0050.AnswerChoice}
  The four answer labels printed beside the source problem
\end{definition}
\begin{proof}
  This is the declared model definition of Answer Choice.
\end{proof}
\begin{definition}[answer Refractive Index]
  \label{def:physics:phyx-mini-0050:phyxminiproblems-problemphyxmini0050-answerrefractiveindex}
  \lean{PhyXMiniProblems.ProblemPhyXMini0050.answerRefractiveIndex}
  \uses{def:physics:phyx-mini-0050:phyxminiproblems-problemphyxmini0050-answerchoice}
  The dimensionless refractive-index value displayed by each answer choice
\end{definition}
\begin{proof}
  This is the declared model definition of answer Refractive Index.
\end{proof}
\begin{definition}[Matches Answer To Nearest Hundredth]
  \label{def:physics:phyx-mini-0050:phyxminiproblems-problemphyxmini0050-matchesanswertonearesthundredth}
  \lean{PhyXMiniProblems.ProblemPhyXMini0050.MatchesAnswerToNearestHundredth}
  \uses{def:physics:phyx-mini-0050:phyxminiproblems-problemphyxmini0050-answerchoice, def:physics:phyx-mini-0050:phyxminiproblems-problemphyxmini0050-answerrefractiveindex}
  Agreement with a refractive index displayed to the nearest hundredth
\end{definition}
\begin{proof}
  This is the declared model definition of Matches Answer To Nearest Hundredth.
\end{proof}
\begin{lemma}[exit interface angles from figure]
  \label{lem:physics:phyx-mini-0050:phyxminiproblems-problemphyxmini0050-exit-interface-angles-from-figure}
  \lean{PhyXMiniProblems.ProblemPhyXMini0050.exit_interface_angles_from_figure}
  \uses{def:physics:phyx-mini-0050:phyxminiproblems-problemphyxmini0050-radiansofdegrees, def:physics:phyx-mini-0050:phyxminiproblems-problemphyxmini0050-prismlasersetup, def:physics:phyx-mini-0050:phyxminiproblems-problemphyxmini0050-anglefrompropagationnormal, def:physics:phyx-mini-0050:phyxminiproblems-problemphyxmini0050-matchesfigurereadouts, def:physics:phyx-mini-0050:phyxminiproblems-problemphyxmini0050-hasdepictedprismgeometry, def:physics:phyx-mini-0050:phyxminiproblems-problemphyxmini0050-hasdepictedraypath}
  The primary-image geometry makes the exit incidence angle 30° and the exit refraction angle 30° + 22.6° = 52.6°, both measured from the exit normal
\end{lemma}
\begin{proof}
  The conclusion follows from the governing relations and figure readouts named in the statement of exit interface angles from figure.
\end{proof}
\begin{lemma}[prism refractive index exact formula]
  \label{lem:physics:phyx-mini-0050:phyxminiproblems-problemphyxmini0050-prism-refractive-index-exact-formula}
  \lean{PhyXMiniProblems.ProblemPhyXMini0050.prism_refractive_index_exact_formula}
  \uses{def:physics:phyx-mini-0050:phyxminiproblems-problemphyxmini0050-dimensionlessreadout, def:physics:phyx-mini-0050:phyxminiproblems-problemphyxmini0050-radiansofdegrees, def:physics:phyx-mini-0050:phyxminiproblems-problemphyxmini0050-prismlasersetup, def:physics:phyx-mini-0050:phyxminiproblems-problemphyxmini0050-matchesfigurereadouts, def:physics:phyx-mini-0050:phyxminiproblems-problemphyxmini0050-hasdepictedprismgeometry, def:physics:phyx-mini-0050:phyxminiproblems-problemphyxmini0050-hasdepictedraypath, def:physics:phyx-mini-0050:phyxminiproblems-problemphyxmini0050-hasphysicalrefractiveindices, def:physics:phyx-mini-0050:phyxminiproblems-problemphyxmini0050-satisfiessnelllawatprismfaces}
  Snell's law at the exit face and the idealized air index determine the exact symbolic prism index before any decimal rounding is performed
\end{lemma}
\begin{proof}
  The conclusion follows from the governing relations and figure readouts named in the statement of prism refractive index exact formula.
\end{proof}
% --- Archon declaration topology end ---
% --- Archon physics formalization source end ---
